\documentclass{article}
\usepackage{amsmath,amssymb,amsthm}
\usepackage{booktabs}
\usepackage{hyperref}
\usepackage{listings}
\usepackage{xcolor}
\usepackage[margin=1in]{geometry}

\newcommand{\hex}[1]{\texttt{\$#1}}
\newcommand{\opcode}[1]{\texttt{#1}}
\newcommand{\Beff}{B_{\mathrm{eff}}}
\newcommand{\Pbase}{P_{\mathrm{base}}}
\newcommand{\Prep}{P_{\mathrm{rep}}}

\newtheorem{definition}{Definition}
\newtheorem{proposition}{Proposition}
\newtheorem{theorem}{Theorem}
\newtheorem{remark}{Remark}

\lstset{
  basicstyle=\ttfamily\small,
  columns=fullflexible,
  frame=single,
  backgroundcolor=\color{gray!10},
}

\title{Probability of Spontaneous Self-Replication\\in a 6502 Cellular Automaton}
\author{}
\date{March 2026}

\begin{document}
\maketitle

\begin{abstract}
We compute the probability that a uniformly random 1024-byte program in a
6502-based cellular automaton constitutes a functional self-replicator.
The strict minimal replicator requires 9 specific bytes, giving a base
probability of $2^{-72}$. We systematically relax the requirements by
cataloging functionally equivalent instruction sequences: alternative
target cells, alternative loop mechanisms, index register variants,
and NOP-insertion tolerance. These relaxations reduce the effective
information content to approximately 43 bits, yielding a probability
of roughly $2^{-43}$. We discuss implications for ``primordial soup''
experiments on random-initialized boards.
\end{abstract}

\tableofcontents

%----------------------------------------------------------------------
\section{System Model}
\label{sec:model}

We consider the 6502life cellular automaton~\cite{6502life}.
Each cell contains $M = 1024$ bytes of RAM.
Execution begins at address \hex{0000} after each scheduler interrupt.
The cell's own memory occupies addresses \hex{0000}--\hex{03FF}
(cell~0 in the memory-mapped neighborhood).
The 48 neighbor cells (in a 7$\times$7 grid minus the origin) are
mapped at addresses \hex{0400}--\hex{BFFF}, with cell~$c$ occupying
addresses $\hex{0400}\cdot c$ through $\hex{0400}\cdot c + \hex{03FF}$
for $c = 1, \ldots, 48$.

The scheduler randomly selects a cell, assigns it a random orientation
(one of four 90$^\circ$ rotations), and resumes execution at the cell's
program counter. After a Poisson-distributed number of cycles, the
scheduler interrupts the cell, saves its registers to fixed locations
in page~0, and moves to the next cell.

\begin{definition}[Byte-level self-replicator]
\label{def:replicator}
A program starting at address \hex{0000} is a \emph{self-replicator}
if, when executed from a state with the X register initialized to some
value $x_0$, it eventually writes a complete copy of its own page~0
(bytes \hex{00}--\hex{FF}) to the page~0 of some neighbor cell~$c$
($1 \le c \le 48$), such that the copy is byte-identical to the original
at least at all positions occupied by the replicator's own code.
\end{definition}

\begin{remark}
This definition is deliberately permissive: the replicator need only
preserve its own code bytes in the copy. Other bytes in page~0 (stack,
data areas) may differ. In practice, the minimal replicators we analyze
copy the entire 256-byte page~0, which is stronger than necessary.
\end{remark}

%----------------------------------------------------------------------
\section{The Minimal Replicator}
\label{sec:minimal}

\subsection{JMP variant (9 bytes)}

The shortest known self-replicator that uses an explicit jump is:

\begin{lstlisting}
      LDA $00,X    ; B5 00     load byte X from own page 0
      STA $0400,X  ; 9D 00 04  store to neighbor cell 1, page 0
      INX          ; E8        increment index
      JMP $0000    ; 4C 00 00  loop
\end{lstlisting}

The byte sequence is: \texttt{B5 00 9D 00 04 E8 4C 00 00}.

This program copies byte~$X$ from its own page~0 to offset~$X$ in
cell~1's page~0, increments $X$, and loops. Since $X$ is an 8-bit
register that wraps from \hex{FF} to \hex{00}, after 256 iterations
every byte of page~0 has been copied, including the replicator's own
code. The copy is therefore byte-identical to the original.

\subsection{BRK variant (7 bytes)}

The opcode \opcode{BRK} (\hex{00}) with operand \hex{00} resets the
program counter to \hex{0000} and yields to the scheduler. If the cell
has enough remaining cycles, execution resumes at \hex{0000}. In
zero-initialized RAM, the byte following \opcode{INX} is already
\hex{00}, so:

\begin{lstlisting}
      LDA $00,X    ; B5 00
      STA $0400,X  ; 9D 00 04
      INX          ; E8
      BRK 0        ; 00 00  (second 00 is BRK operand, already zero)
\end{lstlisting}

Byte sequence: \texttt{B5 00 9D 00 04 E8 00}.
The trailing \hex{00} (BRK operand) is ``free'' in zero-initialized RAM.
The 6502 BRK instruction always reads one operand byte; here it reads
byte~7 which is zero.

\begin{remark}
The BRK variant has a subtlety: \opcode{BRK} triggers the scheduler,
which re-randomizes the cell's orientation and may reassign the neighbor
mapping. However, the copy loop is stateless (X wraps modulo 256), so
after each scheduler re-entry the copy continues from whatever X value
was saved.  The destination cell changes between interrupts, so the
replicator distributes partial copies across multiple neighbors. For
complete self-replication, the replicator must execute 256 consecutive
iterations to a single target without interruption.
The expected number of cycles for 256 iterations of the 4-instruction
loop is approximately $256 \times 10 = 2560$ cycles. Whether this
completes before the Poisson-distributed interrupt depends on the
scheduler's mean quantum.
\end{remark}

\subsection{Branch variant (8 bytes)}

A two-byte branch instruction can replace the three-byte \opcode{JMP}:

\begin{lstlisting}
      LDA $00,X    ; B5 00
      STA $0400,X  ; 9D 00 04
      INX          ; E8
      BNE $F8      ; D0 F8    branch back to offset 0 if Z=0
\end{lstlisting}

Byte sequence: \texttt{B5 00 9D 00 04 E8 D0 F8} (8 bytes).

The branch offset is computed as: from byte~6, \opcode{BNE} is 2 bytes,
so $\mathrm{PC} = 8$ after fetch. Target is byte~0.
Offset $= 0 - 8 = -8$, encoded as the two's complement byte \hex{F8}.

After \opcode{INX}, the zero flag $Z$ is clear unless $X$ wrapped to
zero. So the branch is taken for 255 of every 256 iterations. On the
256th iteration ($X = 0$), the branch falls through, and execution
continues with whatever bytes follow. In zero-initialized RAM, the
next byte is \hex{00} = \opcode{BRK}, which resets PC to 0 and
resumes. So the branch variant is a complete replicator, with a brief
one-time \opcode{BRK} detour every 256 iterations.

%----------------------------------------------------------------------
\section{Base Probability}
\label{sec:base}

\begin{proposition}[Strict base probability]
If all 1024 bytes are drawn independently and uniformly from
$\{0, 1, \ldots, 255\}$, the probability that the first $N$ bytes
match a specific $N$-byte replicator pattern is
\[
   \Pbase = 256^{-N} = 2^{-8N}.
\]
\end{proposition}

\begin{center}
\begin{tabular}{lcc}
\toprule
Variant & $N$ (bytes) & $\Pbase$ \\
\midrule
JMP  & 9 & $2^{-72}$ \\
Branch & 8 & $2^{-64}$ \\
BRK  & 7 & $2^{-56}$ \\
BRK (free trailing zero) & 6 & $2^{-48}$ \\
\bottomrule
\end{tabular}
\end{center}

For the BRK variant, if we only count the 6 ``non-free'' bytes
(assuming the 7th byte is zero in the initial state), then
$\Pbase = 2^{-48}$.

%----------------------------------------------------------------------
\section{Relaxation Analysis: Byte-by-Byte Constraints}
\label{sec:relaxation}

We now analyze each byte position to determine how many of the 256
possible values yield a functionally equivalent replicator.

\subsection{Byte 0: Load opcode}

The load instruction must read a byte from the cell's own page~0 using
an index register. The valid opcodes are:

\begin{center}
\begin{tabular}{llll}
\toprule
Opcode & Mnemonic & Addressing Mode & Register used \\
\midrule
\hex{B5} & \opcode{LDA zp,X} & zero page, X-indexed & X \\
\hex{B6} & \opcode{LDX zp,Y} & zero page, Y-indexed & Y$^\dagger$ \\
\bottomrule
\end{tabular}
\end{center}

$^\dagger$\opcode{LDX zp,Y} loads into the X register, but is indexed
by Y. This requires the corresponding store and increment to also use Y
(or X), which we will account for in those positions.

If we use \opcode{LDA zp,X}, then the store must be \opcode{STA abs,X}
and the increment must be \opcode{INX}. If we use \opcode{LDX zp,Y},
then the store must be a different instruction (since there is no
\opcode{STX abs,Y} on the 6502). We could use
\opcode{LDA zp,Y} (\hex{B6}... wait, \hex{B6} is \opcode{LDX zp,Y}).

Let us be more systematic. The 6502 zero-page indexed addressing modes
for loads are:
\begin{itemize}
\item \hex{B5}: \opcode{LDA zp,X} --- loads A, indexed by X
\item \hex{B6}: \opcode{LDX zp,Y} --- loads X, indexed by Y
\item \hex{B4}: \opcode{LDY zp,X} --- loads Y, indexed by X
\end{itemize}

For a replicator, we need a load--store pair that reads from page~0 and
writes to an absolute indexed address. The absolute indexed stores are:
\begin{itemize}
\item \hex{9D}: \opcode{STA abs,X} --- stores A, indexed by X
\item \hex{99}: \opcode{STA abs,Y} --- stores A, indexed by Y
\end{itemize}
There is no \opcode{STX abs,Y} or \opcode{STY abs,X} in absolute
indexed mode on the 6502.

Therefore the viable load--store--increment combinations are:

\begin{center}
\begin{tabular}{llll}
\toprule
Load & Store & Increment & Index register \\
\midrule
\opcode{LDA zp,X} (\hex{B5}) & \opcode{STA abs,X} (\hex{9D}) & \opcode{INX} (\hex{E8}) or \opcode{DEX} (\hex{CA}) & X \\
\opcode{LDA zp,Y}$^*$ & \opcode{STA abs,Y} (\hex{99}) & \opcode{INY} (\hex{C8}) or \opcode{DEY} (\hex{88}) & Y \\
\bottomrule
\end{tabular}
\end{center}

$^*$There is no official \opcode{LDA zp,Y} instruction on the 6502.
The zero-page Y-indexed addressing mode exists only for \opcode{LDX}
and \opcode{STX}. To load A from a Y-indexed zero-page address, one
would need indirect indexed addressing (\opcode{LDA (zp),Y}, opcode
\hex{B1}), which requires a pointer in zero page and targets a
different address range.

\textbf{Revised analysis:} The only working combination using official
opcodes and absolute indexed addressing is the X-register variant:
\begin{itemize}
\item Load: \hex{B5} (\opcode{LDA zp,X}) --- 1 valid value
\item Store: \hex{9D} (\opcode{STA abs,X}) --- 1 valid value
\item Increment: \hex{E8} (\opcode{INX}) or \hex{CA} (\opcode{DEX}) --- 2 valid values
\end{itemize}

However, we can also construct a replicator using indirect indexed
addressing:

\begin{lstlisting}
      LDA ($00),Y  ; B1 00     load via pointer at $00,$01
      STA ($02),Y  ; 91 02     store via pointer at $02,$03
      INY          ; C8
      BNE $FA      ; D0 FA     branch back
\end{lstlisting}

This requires the zero-page pointers to be set up correctly:
\hex{00},\hex{01} = \hex{00},\hex{00} (pointer to \hex{0000}) and
\hex{02},\hex{03} = \hex{00},\hex{04} (pointer to \hex{0400}). But
the replicator's own code occupies bytes 0--7, which conflicts with
the required pointer values at bytes 0--3. This variant is therefore
\textbf{not self-consistent} and we exclude it.

\textbf{Conclusion for byte 0:} Exactly 1 valid value (\hex{B5}).
Constraint: 8 bits.

\subsection{Byte 1: Zero-page base address}

The operand of \opcode{LDA zp,X} specifies the base address $b$.
The effective address is $(b + X) \bmod 256$. Source byte read when
index register equals $X$ is memory byte $(b + X) \bmod 256$.

The store instruction \opcode{STA \$XX00,X} writes to absolute address
$\hex{XX00} + X$, i.e., offset $X$ in the target cell's page~0.

Therefore, destination byte at offset $k$ receives source byte
$(b + k) \bmod 256$. This is a \emph{rotation} of source page~0 by $b$
positions.

For the copy to be self-replicating, the destination's bytes at the
replicator code positions must match the source's bytes at those same
positions. Let the replicator occupy bytes $0, 1, \ldots, L-1$ and let
bytes $L, L+1, \ldots, 255$ be ``background'' (zero in a fresh cell).

Destination byte $k$ = source byte $(k + b) \bmod 256$.
For $k = 0, 1, \ldots, L-1$: we need source byte $(k+b) \bmod 256$ =
source byte $k$ (the replicator's $k$-th byte).

If $b = 0$, this is trivially satisfied.

If $b \neq 0$, then $(k + b) \bmod 256 \neq k$, and source byte
$(k+b) \bmod 256$ is generically a background byte (zero), not the
replicator byte at position $k$. The only exception would be if the
replicator code happened to be periodic with period dividing $b$, which
our replicator is not (it contains the specific bytes \texttt{B5 00 9D
00 04 E8 ...}, which are not periodic).

\textbf{Conclusion for byte 1:} Must be \hex{00}. Constraint: 8 bits.

\subsection{Byte 2: Store opcode}

Must be \hex{9D} (\opcode{STA abs,X}).
Constraint: 8 bits.

\subsection{Byte 3: Low byte of destination address}

The store instruction \opcode{STA \$HH{LL},X} writes to address
$\hex{HHLL} + X$. The low byte LL determines the base offset within
the target cell's page.

By the same rotation argument as for byte~1: destination byte at
offset $k$ receives the value of the source byte loaded when $X = k$.
The address written is $\hex{HH}\cdot 256 + \mathrm{LL} + X$.
Cell~$c$ occupies addresses $1024c$ through $1024c + 1023$.
The offset within cell~$c$ is $(\mathrm{LL} + X) \bmod 1024$.

For page~0 of cell~$c$ (offsets 0--255): the write hits offset
$(\mathrm{LL} + X) \bmod 256$ of page~0, provided $\hex{HH}$ selects
page~0 of cell~$c$ (which requires $\hex{HH} = 4c$).
If $\mathrm{LL} = 0$, write offset equals $X$, and the copy is
unrotated. If $\mathrm{LL} \neq 0$, the copy is rotated by LL, and the
same analysis as byte~1 applies: the replicator is not self-consistent.

\textbf{Conclusion for byte 3:} Must be \hex{00}. Constraint: 8 bits.

\subsection{Byte 4: High byte of destination address}
\label{sec:byte4}

The high byte $\hex{HH}$ must equal $4c$ for some valid neighbor
cell $c \in \{1, 2, \ldots, 48\}$, so that the store targets page~0
of a neighbor cell.

The valid values are $\{4, 8, 12, \ldots, 192\}$, i.e., the set
$\{4c : 1 \le c \le 48\}$. These are 48 distinct values, all multiples
of~4 in the range $[4, 192]$.

The probability that a uniform random byte is one of these 48 values is
$48/256 = 3/16$.

\textbf{Conclusion for byte 4:} 48 valid values out of 256.
Constraint: $\log_2(256/48) = \log_2(16/3) \approx 2.415$ bits.

\subsection{Byte 5: Increment/decrement opcode}

The index register X must change by $\pm 1$ each iteration so that all
256 byte offsets are visited.

\begin{itemize}
\item \hex{E8}: \opcode{INX} --- increments X. Copies bytes in order
  $0, 1, 2, \ldots, 255$ (starting from initial $X$).
\item \hex{CA}: \opcode{DEX} --- decrements X. Copies bytes in order
  $0, 255, 254, \ldots, 1$. All 256 bytes are still copied.
\end{itemize}

No other single-byte instruction changes X by exactly $\pm 1$ without
side effects that would corrupt the replicator.

Valid values: 2 out of 256. Constraint:
$\log_2(256/2) = 7$ bits.

\subsection{Bytes 6+: Loop instruction}
\label{sec:loop}

We analyze three loop mechanisms.

\subsubsection{JMP variant (bytes 6--8)}

\opcode{JMP \$0000} is encoded as \texttt{4C 00 00} (3 bytes).
The operand must be exactly \hex{0000} to return to the start.
Constraint: $8 + 8 + 8 = 24$ bits (3 bytes).

No alternative: \opcode{JMP} is the only 6502 instruction that performs
an unconditional absolute jump. (\opcode{JMP (indirect)} at \hex{6C}
requires a pointer, adding complexity.)

\subsubsection{Branch variant (bytes 6--7)}

A conditional branch instruction is 2 bytes: opcode plus signed offset.
The offset must be \hex{F8} ($= -8$ signed) to branch back to byte~0
from byte~6 (since PC after fetch = $6 + 2 = 8$, and $8 + (-8) = 0$).

The branch is taken if its condition is met. After \opcode{INX} (or
\opcode{DEX}), the processor flags are:
\begin{itemize}
\item $Z = 1$ iff $X = 0$ (wrapped). Otherwise $Z = 0$.
\item $N = 1$ iff bit~7 of $X$ is set (i.e., $X \in [128, 255]$).
  Otherwise $N = 0$.
\item $C$ is \emph{unchanged} by \opcode{INX}/\opcode{DEX}.
\item $V$ is \emph{unchanged}.
\end{itemize}

The branches and their reliability:

\begin{center}
\begin{tabular}{lllc}
\toprule
Opcode & Mnemonic & Condition & Iterations taken (of 256) \\
\midrule
\hex{D0} & \opcode{BNE} & $Z = 0$ & 255 \\
\hex{F0} & \opcode{BEQ} & $Z = 1$ & 1 \\
\hex{10} & \opcode{BPL} & $N = 0$ & 128 \\
\hex{30} & \opcode{BMI} & $N = 1$ & 128 \\
\hex{90} & \opcode{BCC} & $C = 0$ & depends on entry state \\
\hex{B0} & \opcode{BCS} & $C = 1$ & depends on entry state \\
\hex{50} & \opcode{BVC} & $V = 0$ & depends on entry state \\
\hex{70} & \opcode{BVS} & $V = 1$ & depends on entry state \\
\bottomrule
\end{tabular}
\end{center}

For a \emph{complete} replicator (copies all 256 bytes without
interruption), we need the branch taken for 255 consecutive iterations.
Only \opcode{BNE} achieves this: it fails on exactly the one iteration
where $X$ wraps to 0, at which point all 256 bytes have been copied.

\opcode{BPL} and \opcode{BMI} fail on half the iterations, so the loop
copies at most a few bytes before falling through. \opcode{BEQ} almost
never branches. The carry- and overflow-dependent branches (\opcode{BCC},
\opcode{BCS}, \opcode{BVC}, \opcode{BVS}) depend on prior state:
\opcode{INX}/\opcode{DEX} does not modify C or V, so if C was clear on
entry, \opcode{BCC} always branches (all 256 iterations, including the
wrap), and if C was set, \opcode{BCC} never branches.

In a randomly initialized cell, the carry flag has probability 1/2 of
being in either state. Therefore:
\begin{itemize}
\item \opcode{BCC} (\hex{90}): works if $C=0$ on entry (probability 1/2).
\item \opcode{BCS} (\hex{B0}): works if $C=1$ on entry (probability 1/2).
\item \opcode{BVC} (\hex{50}): works if $V=0$ on entry (probability 1/2).
\item \opcode{BVS} (\hex{70}): works if $V=1$ on entry (probability 1/2).
\end{itemize}

For \opcode{BNE}: the branch is not taken when $X$ wraps to 0, but by
then all 256 bytes have been copied, and in zero-initialized RAM the
fall-through byte is \hex{00} (\opcode{BRK}), which resets PC to 0.
So the replicator is functional.

For the flag-dependent branches, if the flag is in the correct state,
the branch is taken on all 256 iterations (the loop never terminates
naturally, but the scheduler eventually interrupts). The copy may or
may not complete within one scheduler quantum, but statistically it
will complete over multiple quanta since $X$ is saved and restored.

\textbf{Counting valid branch opcodes:}
\begin{itemize}
\item Unconditionally reliable: \opcode{BNE} (\hex{D0}) --- 1 opcode.
\item Reliable with probability 1/2: \opcode{BCC} (\hex{90}),
  \opcode{BCS} (\hex{B0}), \opcode{BVC} (\hex{50}),
  \opcode{BVS} (\hex{70}) --- 4 opcodes.
\item Unreliable: \opcode{BPL} (\hex{10}), \opcode{BMI} (\hex{30}),
  \opcode{BEQ} (\hex{F0}) --- excluded.
\end{itemize}

For a conservative count, we take only \opcode{BNE}: 1 valid opcode.
For a generous count including the flag-dependent branches with their
1/2 probability weight:
\[
  \text{effective count} = 1 + 4 \times \tfrac{1}{2} = 3.
\]

Byte~6 constraint (generous): $\log_2(256/3) \approx 6.415$ bits.\\
Byte~7 constraint: must be \hex{F8} exactly. 8 bits.

Total for branch loop (generous): $\approx 14.415$ bits over 2 bytes.

\subsubsection{BRK variant (bytes 6+)}
\label{sec:brk}

If byte~6 is \hex{00} (\opcode{BRK}), the scheduler resets PC to~0.
The BRK instruction reads one operand byte (byte~7); in
zero-initialized RAM, this is already \hex{00}.

The critical question is whether byte~6 must be exactly \hex{00}, or
whether it is ``free'' because it falls in the zero-initialized
background.

\textbf{Case 1: Explicit BRK.} We require byte~6 = \hex{00}. Since
the remaining bytes (7+) are unconstrained (they are zero by
initialization, and the BRK operand is ignored beyond being read),
we have 1 constrained byte.
Constraint: 8 bits for byte~6.

\textbf{Case 2: Implicit BRK.} If we assume zero initialization, then
any program of length $\le 6$ bytes is followed by \hex{00} bytes.
A 6-byte replicator (bytes 0--5: \texttt{B5 00 9D 00 04 E8}) is
followed by \hex{00} at byte~6 ``for free.''
Constraint: 0 additional bits.

In this paper, we treat both cases. For random initialization (all
bytes uniform), byte~6 must be \hex{00} explicitly (Case~1). For
zero-initialized memory, only 6 bytes are constrained (Case~2).

%----------------------------------------------------------------------
\section{Summary of Byte Constraints}
\label{sec:summary}

We present the effective information content for each variant.

\subsection{JMP variant (9 bytes, random initialization)}

\begin{center}
\begin{tabular}{clccc}
\toprule
Byte & Content & Valid values & Probability & Bits \\
\midrule
0 & \opcode{LDA zp,X} opcode & 1 (/256) & $2^{-8}$ & 8.000 \\
1 & Zero-page base addr & 1 (/256) & $2^{-8}$ & 8.000 \\
2 & \opcode{STA abs,X} opcode & 1 (/256) & $2^{-8}$ & 8.000 \\
3 & Dest addr low byte & 1 (/256) & $2^{-8}$ & 8.000 \\
4 & Dest addr high byte & 48 (/256) & $48/256$ & 2.415 \\
5 & \opcode{INX} or \opcode{DEX} & 2 (/256) & $2/256$ & 7.000 \\
6 & \opcode{JMP} opcode & 1 (/256) & $2^{-8}$ & 8.000 \\
7 & Jump addr low byte & 1 (/256) & $2^{-8}$ & 8.000 \\
8 & Jump addr high byte & 1 (/256) & $2^{-8}$ & 8.000 \\
\midrule
\multicolumn{4}{r}{\textbf{Total}} & \textbf{65.415} \\
\bottomrule
\end{tabular}
\end{center}

\[
  \Prep^{(\mathrm{JMP})} = \frac{1 \cdot 1 \cdot 1 \cdot 1 \cdot 48 \cdot 2 \cdot 1 \cdot 1 \cdot 1}{256^9} = \frac{96}{256^9} = \frac{96}{2^{72}} = \frac{3}{2^{67}} \approx 2^{-65.415}.
\]

\subsection{Branch variant (8 bytes, random initialization)}

\begin{center}
\begin{tabular}{clccc}
\toprule
Byte & Content & Valid values & Probability & Bits \\
\midrule
0 & \opcode{LDA zp,X} opcode & 1 (/256) & $2^{-8}$ & 8.000 \\
1 & Zero-page base addr & 1 (/256) & $2^{-8}$ & 8.000 \\
2 & \opcode{STA abs,X} opcode & 1 (/256) & $2^{-8}$ & 8.000 \\
3 & Dest addr low byte & 1 (/256) & $2^{-8}$ & 8.000 \\
4 & Dest addr high byte & 48 (/256) & $48/256$ & 2.415 \\
5 & \opcode{INX} or \opcode{DEX} & 2 (/256) & $2/256$ & 7.000 \\
6 & Branch opcode & 1--3 (/256) & $1/256$--$3/256$ & $6.415$--$8$ \\
7 & Branch offset & 1 (/256) & $2^{-8}$ & 8.000 \\
\midrule
\multicolumn{4}{r}{\textbf{Total (conservative, BNE only)}} & \textbf{57.415} \\
\multicolumn{4}{r}{\textbf{Total (generous, BNE + flag-dependent)}} & \textbf{55.830} \\
\bottomrule
\end{tabular}
\end{center}

Conservative (BNE only):
\[
  \Prep^{(\mathrm{BNE})} = \frac{1 \cdot 1 \cdot 1 \cdot 1 \cdot 48 \cdot 2 \cdot 1 \cdot 1}{256^8}
  = \frac{96}{2^{64}} = \frac{3}{2^{59}} \approx 2^{-57.415}.
\]

Generous (effective 3 branch opcodes):
\[
  \Prep^{(\mathrm{branch})} = \frac{1 \cdot 1 \cdot 1 \cdot 1 \cdot 48 \cdot 2 \cdot 3 \cdot 1}{256^8}
  = \frac{288}{2^{64}} = \frac{9}{2^{59}} \approx 2^{-55.830}.
\]

\subsection{BRK variant (random initialization)}

\begin{center}
\begin{tabular}{clccc}
\toprule
Byte & Content & Valid values & Probability & Bits \\
\midrule
0 & \opcode{LDA zp,X} opcode & 1 (/256) & $2^{-8}$ & 8.000 \\
1 & Zero-page base addr & 1 (/256) & $2^{-8}$ & 8.000 \\
2 & \opcode{STA abs,X} opcode & 1 (/256) & $2^{-8}$ & 8.000 \\
3 & Dest addr low byte & 1 (/256) & $2^{-8}$ & 8.000 \\
4 & Dest addr high byte & 48 (/256) & $48/256$ & 2.415 \\
5 & \opcode{INX} or \opcode{DEX} & 2 (/256) & $2/256$ & 7.000 \\
6 & \opcode{BRK} opcode & 1 (/256) & $2^{-8}$ & 8.000 \\
\midrule
\multicolumn{4}{r}{\textbf{Total (7 bytes, random init)}} & \textbf{49.415} \\
\bottomrule
\end{tabular}
\end{center}

\[
  \Prep^{(\mathrm{BRK,7})} = \frac{96}{256^7} = \frac{96}{2^{56}} = \frac{3}{2^{51}} \approx 2^{-49.415}.
\]

With zero initialization (byte~6 is free):
\[
  \Prep^{(\mathrm{BRK,6})} = \frac{96}{256^6} = \frac{96}{2^{48}} = \frac{3}{2^{43}} \approx 2^{-41.415}.
\]

\subsection{Combined probability (union over all variants)}

A random byte sequence is a valid replicator if it matches \emph{any}
of the variants. Since the variants have disjoint byte patterns at
position 6 (\hex{4C} for JMP, \hex{D0}/\hex{90}/\hex{B0}/\hex{50}/\hex{70}
for branches, \hex{00} for BRK), their probabilities are additive up to
the shared prefix:

\begin{align}
  \Prep^{(\mathrm{total})} &= \Prep^{(\mathrm{JMP})} + \Prep^{(\mathrm{branch})} + \Prep^{(\mathrm{BRK,7})} \nonumber \\
  &= \frac{96}{256^9} + \frac{288}{256^8} + \frac{96}{256^7} \nonumber \\
  &= \frac{96}{2^{72}} + \frac{288}{2^{64}} + \frac{96}{2^{56}} \label{eq:total}
\end{align}

The BRK term dominates:
\[
  \Prep^{(\mathrm{total})} \approx \frac{96}{2^{56}} = \frac{3}{2^{51}} \approx 2^{-49.415}.
\]

The JMP and branch terms are negligible by comparison (smaller by
factors of $2^{-16}$ and $2^{-8}$ respectively).

%----------------------------------------------------------------------
\section{NOP Insertion Tolerance}
\label{sec:nop}

\subsection{Setup}

The replicator's four instructions (load, store, increment, loop) are
independent in the sense that inserting ``no-operation'' instructions
between them does not change the computation. A NOP does not modify
registers A, X, Y, or the status flags relevant to the branch condition.

We allow NOP insertions at $K = 4$ positions:
\begin{enumerate}
\item Before the load (leading NOPs, $j_0$ bytes)
\item Between load and store ($j_1$ bytes)
\item Between store and increment ($j_2$ bytes)
\item Between increment and loop ($j_3$ bytes)
\end{enumerate}

Let $J = j_0 + j_1 + j_2 + j_3$ be the total number of NOP bytes
inserted. The resulting program has length $N_0 + J$ bytes, where
$N_0$ is the base replicator length.

\subsection{Counting NOP opcodes}

The official 6502 NOP is \hex{EA} (1 byte, 2 cycles). The 6502 also
has several undocumented single-byte NOPs that are functionally
equivalent (no register or flag side effects):
\[
  \hex{1A},\ \hex{3A},\ \hex{5A},\ \hex{7A},\ \hex{DA},\ \hex{FA}.
\]
This gives $n_{\mathrm{NOP}} = 7$ single-byte NOP opcodes.

There are also multi-byte undocumented NOPs (2-byte and 3-byte), but
these complicate the analysis significantly because their operand bytes
add additional degrees of freedom. We defer multi-byte NOPs to
Section~\ref{sec:multibyte-nop}.

\subsection{Generating function for NOP insertions}

For each insertion point, the number of single-byte NOP bytes inserted
is $j_i \ge 0$, and each such byte must be one of $n_{\mathrm{NOP}}$
values. The probability of $j$ consecutive NOP bytes at a single
insertion point is
\[
  \left(\frac{n_{\mathrm{NOP}}}{256}\right)^j.
\]

The total probability contribution from NOP-padded replicators is:
\begin{align}
  \Prep^{(\mathrm{NOP})} &= P_{\mathrm{core}} \cdot \sum_{j_0=0}^{\infty} \sum_{j_1=0}^{\infty} \sum_{j_2=0}^{\infty} \sum_{j_3=0}^{\infty} \left(\frac{n_{\mathrm{NOP}}}{256}\right)^{j_0+j_1+j_2+j_3} \cdot A(j_0) \label{eq:nop-sum}
\end{align}
where $P_{\mathrm{core}}$ is the probability of the core instruction
bytes (excluding the loop-back), and $A(j_0)$ is a correction factor
for the loop instruction's operand, which depends on the total program
length and hence on the NOP counts.

\textbf{Branch variant correction.}
For the branch variant, the branch offset depends on the total program
length $L = N_0 + J$. From byte position $L - 2$ (the branch opcode),
the required offset to reach byte~0 is $-(L)$, encoded as a signed byte.
This is valid only if $L \le 128$ (the maximum backward branch range
is $-128$). For $L \le 128$, the offset byte is uniquely determined
($256 - L$). So $A(j_0) = 1$ for all valid $J$, and the branch
offset is always exactly 1 value.

However, there is a subtlety for leading NOPs ($j_0 > 0$): the loop
branches back to byte~0, which is a NOP, not the load instruction.
The replicator still works because the NOPs are re-executed each
iteration. The branch offset must target byte~0, not the first
non-NOP instruction.

For the branch variant with leading NOPs, the branch offset is
$-L = -(N_0 + J)$. This is valid as a signed byte when $L \le 128$,
i.e., $J \le 128 - N_0 = 120$ (for $N_0 = 8$). Each valid offset is
one specific byte value. So $A(j_0) = 1$ and the constraint on the
offset byte remains 8 bits regardless of $J$.

\textbf{JMP variant correction.}
The \opcode{JMP \$0000} operand is always \hex{00 00} regardless of
program length. So $A(j_0) = 1$ trivially.

\textbf{BRK variant correction.}
\opcode{BRK} resets PC to 0 regardless of where it appears.
$A(j_0) = 1$ trivially.

Therefore, in all cases $A(j_0) = 1$, and Equation~\eqref{eq:nop-sum}
simplifies:
\begin{align}
  \Prep^{(\mathrm{NOP})} &= P_{\mathrm{core}} \cdot \left(\sum_{j=0}^{\infty} \left(\frac{n_{\mathrm{NOP}}}{256}\right)^j\right)^K \nonumber \\
  &= P_{\mathrm{core}} \cdot \left(\frac{1}{1 - n_{\mathrm{NOP}}/256}\right)^K \nonumber \\
  &= P_{\mathrm{core}} \cdot \left(\frac{256}{256 - n_{\mathrm{NOP}}}\right)^K. \label{eq:nop-factor}
\end{align}

\begin{proposition}[NOP correction factor]
\label{prop:nop}
With $n_{\mathrm{NOP}} = 7$ single-byte NOP opcodes and $K = 4$
insertion points, the NOP-tolerance multiplicative factor is
\[
  F_{\mathrm{NOP}} = \left(\frac{256}{249}\right)^4 \approx 1.1174.
\]
This increases the replicator probability by a factor of $\approx 1.12$,
reducing the effective information content by
\[
  \Delta B = \log_2 F_{\mathrm{NOP}} = 4 \log_2\!\left(\frac{256}{249}\right) \approx 0.160 \text{ bits}.
\]
\end{proposition}

\begin{proof}
The geometric series $\sum_{j=0}^{\infty} x^j = (1-x)^{-1}$ for
$|x| < 1$. With $x = n_{\mathrm{NOP}}/256 = 7/256 \approx 0.02734$,
the series converges and equals $256/249$. Raising to the $K=4$th power:
\[
  F_{\mathrm{NOP}} = \left(\frac{256}{249}\right)^4.
\]
Numerically: $256/249 \approx 1.02811$, so
$F_{\mathrm{NOP}} \approx 1.02811^4 \approx 1.1174$.
And $\log_2(1.1174) \approx 0.160$.
\end{proof}

\subsection{Multi-byte NOP instructions}
\label{sec:multibyte-nop}

The 6502 has undocumented 2-byte NOPs (opcode + 1 operand byte) and
3-byte NOPs (opcode + 2 operand bytes) that have no side effects.
These provide additional NOP-equivalent sequences:

\begin{itemize}
\item \textbf{2-byte NOPs:} Opcodes \hex{80}, \hex{82}, \hex{C2},
  \hex{E2}, \hex{04}, \hex{14}, \hex{34}, \hex{44}, \hex{54},
  \hex{64}, \hex{74}, \hex{D4}, \hex{F4}. Each takes 1 operand byte
  that can be anything (256 values). This gives $n_2 \approx 13$
  distinct 2-byte NOP opcodes.
  Probability per 2-byte NOP: $(n_2/256) \cdot 1 = n_2/256$.

\item \textbf{3-byte NOPs:} Opcodes \hex{0C}, \hex{1C}, \hex{3C},
  \hex{5C}, \hex{7C}, \hex{DC}, \hex{FC}. Each takes 2 operand bytes.
  This gives $n_3 = 7$ distinct 3-byte NOP opcodes.
  Probability per 3-byte NOP: $(n_3/256) \cdot 1 \cdot 1 = n_3/256$.
\end{itemize}

Wait --- a 2-byte NOP inserted at a gap adds 2 bytes to the program
length, with probability $(n_2/256) \cdot (256/256) = n_2/256$ for the
pair (the operand byte is unconstrained, contributing a factor of
$256/256 = 1$). Similarly, a 3-byte NOP adds 3 bytes with probability
$n_3/256$.

Actually, we must be more careful. Each NOP ``slot'' can hold a
sequence of NOP instructions of various lengths. Let $p_k$ be the
probability that a specific $k$-byte NOP instruction appears, consuming
$k$ bytes from the random stream:
\begin{align*}
  p_1 &= n_1 / 256 = 7/256, \\
  p_2 &= n_2 / 256 = 13/256, \quad \text{(operand byte is free)} \\
  p_3 &= n_3 / 256 = 7/256. \quad \text{(operand bytes are free)}
\end{align*}

Note that for multi-byte NOPs, we ``consume'' additional random bytes
from the stream, but those bytes are unconstrained (any value is a valid
operand). The probability cost is only the opcode byte.

At each insertion point, the content is a sequence of NOP instructions.
Let $G(z)$ be the generating function where $z$ tracks the number of
random bytes consumed. A single NOP instruction consuming $k$ bytes
contributes $p_k z^k$ to the generating function. The generating
function for zero or more NOPs at one insertion point is:
\[
  H(z) = \sum_{m=0}^{\infty} \left(\sum_{k=1}^{3} p_k z^k \right)^m
  = \frac{1}{1 - (p_1 z + p_2 z^2 + p_3 z^3)}.
\]

The total probability of a NOP-padded replicator is obtained by
evaluating $H(z)$ at $z = 1/256$ for each consumed byte... but wait,
this is not quite right. The generating function approach needs to
account for the fact that each random byte costs a factor of $1/256$
in probability.

Let us re-derive. Suppose a NOP insertion point consumes $L$ bytes
total. The probability that $L$ specific random bytes form a valid
sequence of NOP instructions is:
\[
  Q(L) = \sum_{\text{NOP sequences of length } L} \prod_{\text{instructions}} \frac{1}{256^{k_i}}
\]
where $k_i$ is the length of the $i$th NOP instruction.

For a $k$-byte NOP instruction, the opcode must be one of $n_k$ values
(probability $n_k/256$) and the remaining $k-1$ operand bytes are free
(probability $1$ each). So the probability cost per instruction is
$n_k / 256$, independent of instruction length. This is because the
$k-1$ operand bytes are unconstrained --- they consume $k-1$ bytes
from the random stream but impose no constraint.

Wait, that is wrong. Each byte in the random 1024-byte cell is drawn
uniformly. A $k$-byte NOP instruction ``uses up'' $k$ bytes. The first
byte must be a valid NOP opcode ($n_k / 256$ probability). The remaining
$k-1$ bytes can be anything ($1$ probability each). But these $k-1$
bytes are still random bytes that could have been something else. The
total probability that bytes $i, i+1, \ldots, i+k-1$ form a specific
$k$-byte NOP is $(n_k / 256) \cdot 1^{k-1} = n_k / 256$.

But the remaining $k-1$ bytes are ``consumed'' --- they cannot be part
of another NOP instruction or the replicator core. This consumption is
already accounted for by the fact that the replicator has a specific
total length. Let $L_{\mathrm{total}}$ be the total program length
(core + all NOP bytes). Then $\Prep = 256^{-L_{\mathrm{total}}} \times$
(number of valid patterns of that length). Actually, since some bytes
within multi-byte NOPs are unconstrained, the probability per byte
is not uniformly $1/256$.

Let me restart with a cleaner framework.

\begin{definition}[NOP insertion probability]
Fix the core replicator bytes. A \emph{NOP-padded replicator of total
length $L$} is a sequence of $L$ bytes that consists of the core
replicator instructions interspersed with NOP instruction sequences
at the $K$ designated insertion points, such that the total byte count
is~$L$. The probability of a random $L$-byte prefix matching such a
pattern is
\[
  P(L) = P_{\mathrm{core}} \cdot \prod_{i=1}^{K} q_i(\ell_i)
\]
where $\ell_i$ is the number of NOP bytes at insertion point $i$
(with $\sum_i \ell_i = L - N_0$), $P_{\mathrm{core}}$ accounts for
the constrained core bytes, and $q_i(\ell)$ is the probability that
$\ell$ consecutive random bytes form a valid sequence of NOP
instructions.
\end{definition}

The function $q(\ell)$ (same for all insertion points) satisfies:
\begin{itemize}
\item $q(0) = 1$ (no NOP bytes).
\item For $\ell \ge 1$: the first byte starts a NOP instruction of
  length $k \in \{1, 2, 3\}$, with $k \le \ell$. The probability
  that the first byte is a valid $k$-byte NOP opcode is $n_k / 256$.
  The next $k-1$ bytes are unconstrained (free). The remaining
  $\ell - k$ bytes must also form a valid NOP sequence.
\end{itemize}

\[
  q(\ell) = \sum_{k=1}^{\min(\ell,3)} \frac{n_k}{256} \cdot q(\ell - k)
  \quad \text{for } \ell \ge 1.
\]

This is a linear recurrence. The generating function is
$Q(z) = \sum_{\ell=0}^{\infty} q(\ell) z^\ell$, satisfying:
\[
  Q(z) = 1 + \left(\frac{n_1}{256} z + \frac{n_2}{256} z^2 + \frac{n_3}{256} z^3\right) Q(z)
\]
so
\begin{equation}
  Q(z) = \frac{1}{1 - \frac{1}{256}(n_1 z + n_2 z^2 + n_3 z^3)}.
  \label{eq:nop-gf}
\end{equation}

The total probability summing over all possible NOP insertion lengths is:
\begin{align}
  \Prep^{(\mathrm{NOP})} &= P_{\mathrm{core}} \cdot \sum_{\ell_1, \ldots, \ell_K \ge 0} \prod_{i=1}^K q(\ell_i) \nonumber \\
  &= P_{\mathrm{core}} \cdot \left(\sum_{\ell=0}^{\infty} q(\ell)\right)^K \nonumber \\
  &= P_{\mathrm{core}} \cdot Q(1)^K. \label{eq:nop-total}
\end{align}

\begin{proposition}[NOP factor with multi-byte NOPs]
\label{prop:nop-multi}
With $n_1 = 7$, $n_2 = 13$, $n_3 = 7$ and $K = 4$:
\[
  Q(1) = \frac{1}{1 - (7 + 13 + 7)/256} = \frac{256}{256 - 27} = \frac{256}{229} \approx 1.1179.
\]
The NOP correction factor is
\[
  F_{\mathrm{NOP}} = Q(1)^4 = \left(\frac{256}{229}\right)^4 \approx 1.5608.
\]
This reduces the effective information content by
\[
  \Delta B = \log_2(1.5608) \approx 0.642 \text{ bits}.
\]
\end{proposition}

\begin{proof}
From Equation~\eqref{eq:nop-gf} with $z=1$:
\[
  Q(1) = \frac{1}{1 - (n_1 + n_2 + n_3)/256} = \frac{256}{256 - 27} = \frac{256}{229}.
\]
Numerically: $256/229 \approx 1.11790$, and $(256/229)^4 \approx 1.5608$.
Then $\log_2(1.5608) \approx 0.6424$.
\end{proof}

\begin{remark}
The multi-byte NOP generating function Equation~\eqref{eq:nop-gf}
correctly accounts for the fact that multi-byte NOP operand bytes are
unconstrained: $q(\ell)$ gives the probability that $\ell$ consecutive
\emph{random} bytes form a valid NOP sequence, where each random byte
independently has probability $1/256$ of being any specific value.
The ``free'' operand bytes contribute factors of $256/256 = 1$ to the
probability, which is why $n_k/256$ (not $n_k/256^k$) appears in the
recurrence.
\end{remark}

\begin{remark}[Truncation effects]
The generating function $Q(1)$ sums over $\ell = 0, 1, 2, \ldots, \infty$.
In practice, the NOP sequence cannot exceed $1024 - N_0$ bytes.
Since $q(\ell)$ decays exponentially as $(27/256)^\ell \approx 0.1055^\ell$,
the tail beyond $\ell = 10$ contributes less than $(0.1055)^{10} \approx
1.5 \times 10^{-10}$ to $Q(1)$. The infinite sum is an excellent
approximation.
\end{remark}

%----------------------------------------------------------------------
\section{Additional Functional Equivalences}
\label{sec:other}

\subsection{Starting position within the cell}

The replicator must start at the entry point, which is byte~0 after
each scheduler interrupt. Leading NOPs (the $j_0$ insertion point) are
the only way to shift the ``real'' code away from byte~0 while
maintaining functionality. This is already accounted for in
Section~\ref{sec:nop}.

\subsection{Alternative copy directions}

Using \opcode{DEX} (\hex{CA}) instead of \opcode{INX} (\hex{E8}) copies
bytes in the order $x_0, x_0 - 1, \ldots, 0, 255, \ldots, x_0 + 1$.
All 256 bytes are still copied. This is already counted (2 valid values
at byte~5).

\subsection{Copy to different cells}

As analyzed in Section~\ref{sec:byte4}, any of the 48 neighbor cells
is a valid target. This is already counted (48 valid values at byte~4).

\subsection{Partial copies and multi-quantum replication}

The analysis above assumes the replicator completes a full 256-byte
copy within a single scheduler quantum. In practice, the replicator
may be interrupted partway through and resume in a later quantum (to a
different, randomly chosen neighbor). This ``scattered copy'' mode is
less reliable but could still result in eventual replication if the
same neighbor is targeted enough times. We do not quantify this here;
our analysis is conservative in requiring a single-quantum complete
copy.

\subsection{Self-modifying or longer replicators}

Longer programs could implement replication through entirely different
mechanisms (e.g., using a subroutine, using indirect addressing, using
self-modifying code). These are vastly more complex and correspondingly
less probable under random initialization. We do not attempt to
enumerate them.

%----------------------------------------------------------------------
\section{Effective Information Content}
\label{sec:beff}

Combining all relaxations from Sections~\ref{sec:relaxation}
and~\ref{sec:nop}, we compute the effective information content
$\Beff$, defined as $\Beff = -\log_2 \Prep$.

\subsection{BRK variant, random initialization (most favorable)}

Core probability (7 bytes):
\[
  P_{\mathrm{core}}^{(\mathrm{BRK,7})} = \frac{1 \cdot 1 \cdot 1 \cdot 1 \cdot 48 \cdot 2 \cdot 1}{256^7} = \frac{96}{2^{56}} = \frac{3}{2^{51}}.
\]

With NOP correction (Proposition~\ref{prop:nop-multi}):
\[
  \Prep^{(\mathrm{BRK,NOP})} = \frac{3}{2^{51}} \times \left(\frac{256}{229}\right)^4 \approx \frac{3}{2^{51}} \times 1.5608.
\]

\[
  \Beff^{(\mathrm{BRK,NOP})} = 51 - \log_2 3 - \log_2 1.5608 \approx 51 - 1.585 - 0.642 \approx 48.773 \text{ bits}.
\]

\subsection{Branch variant, random initialization, generous opcode count}

Core probability (8 bytes, 3 effective branch opcodes):
\[
  P_{\mathrm{core}}^{(\mathrm{branch})} = \frac{1 \cdot 1 \cdot 1 \cdot 1 \cdot 48 \cdot 2 \cdot 3 \cdot 1}{256^8} = \frac{288}{2^{64}} = \frac{9}{2^{59}}.
\]

With NOP correction:
\[
  \Prep^{(\mathrm{branch,NOP})} = \frac{9}{2^{59}} \times \left(\frac{256}{229}\right)^4.
\]

\[
  \Beff^{(\mathrm{branch,NOP})} = 59 - \log_2 9 - \log_2 1.5608 \approx 59 - 3.170 - 0.642 \approx 55.188 \text{ bits}.
\]

\subsection{Union over all variants with NOP correction}

The total probability is dominated by the BRK variant (shortest):
\[
  \Prep^{(\mathrm{total})} \approx \Prep^{(\mathrm{BRK,NOP})} \approx 2^{-48.773}.
\]

\begin{theorem}[Effective information content]
\label{thm:beff}
The effective information content of the minimal 6502 self-replicator,
accounting for alternative target cells (48 neighbors), bidirectional
index stepping (\opcode{INX}/\opcode{DEX}), and NOP insertion tolerance
(7 single-byte + 13 two-byte + 7 three-byte undocumented NOP opcodes
at 4 insertion points), is
\[
  \boxed{\Beff \approx 48.8 \text{ bits}}
\]
for the BRK variant under random initialization, or
$\Beff \approx 40.8$ bits under zero initialization (where the trailing
BRK is free).
\end{theorem}

%----------------------------------------------------------------------
\section{Primordial Soup: Probability of Spontaneous Emergence}
\label{sec:primordial}

\subsection{Expected number of replicators on a random board}

Consider a board of $B \times B$ cells, each with $M = 1024$ bytes
drawn uniformly at random. The replicator must start at byte~0
(the execution entry point), so each cell has exactly one ``trial''
for the replicator pattern.

\begin{proposition}
The expected number of replicators on a $B \times B$ board is
\[
  \mathbb{E}[R] = B^2 \cdot \Prep.
\]
\end{proposition}

For the standard 256$\times$256 board ($B^2 = 65{,}536 = 2^{16}$):
\[
  \mathbb{E}[R] = 2^{16} \cdot 2^{-\Beff}.
\]

\begin{center}
\begin{tabular}{lcc}
\toprule
Scenario & $\Beff$ & $\mathbb{E}[R]$ on $256^2$ board \\
\midrule
BRK, random init, with NOPs & 48.8 & $2^{-32.8} \approx 1.4 \times 10^{-10}$ \\
BRK, zero init, with NOPs & 40.8 & $2^{-24.8} \approx 3.4 \times 10^{-8}$ \\
Branch, random init, with NOPs & 55.2 & $2^{-39.2} \approx 1.7 \times 10^{-12}$ \\
JMP, random init, with NOPs & 64.8 & $2^{-48.8} \approx 2.8 \times 10^{-15}$ \\
\bottomrule
\end{tabular}
\end{center}

In all cases, the expected number of replicators on a standard board is
astronomically small. Even the most favorable scenario (BRK variant,
zero initialization) gives $\mathbb{E}[R] \approx 10^{-8}$.

\subsection{Threshold board size}

For $\mathbb{E}[R] = 1$, we need $B^2 = 2^{\Beff}$, i.e.,
$B = 2^{\Beff/2}$.

\begin{center}
\begin{tabular}{lcc}
\toprule
Scenario & $\Beff$ & Board side $B$ for $\mathbb{E}[R]=1$ \\
\midrule
BRK, zero init, with NOPs & 40.8 & $2^{20.4} \approx 1.4 \times 10^6$ \\
BRK, random init, with NOPs & 48.8 & $2^{24.4} \approx 2.2 \times 10^7$ \\
Branch, random init, with NOPs & 55.2 & $2^{27.6} \approx 2.0 \times 10^8$ \\
JMP, random init, with NOPs & 64.8 & $2^{32.4} \approx 5.5 \times 10^9$ \\
\bottomrule
\end{tabular}
\end{center}

\subsection{Memory requirements}

Each cell requires 1024 bytes = 1 KiB. A board of side $B$ requires
$B^2 \times 1024$ bytes = $B^2$ KiB.

\begin{center}
\begin{tabular}{lcc}
\toprule
Scenario & Board side $B$ & Memory \\
\midrule
BRK, zero init & $1.4 \times 10^6$ & $\approx 1.8$ TiB \\
BRK, random init & $2.2 \times 10^7$ & $\approx 450$ TiB \\
Branch, random init & $2.0 \times 10^8$ & $\approx 37$ PiB \\
JMP, random init & $5.5 \times 10^9$ & $\approx 28$ EiB \\
\bottomrule
\end{tabular}
\end{center}

Even the most favorable scenario requires $\sim$1.8 TiB of memory for
a board where the expected number of spontaneous replicators is~1. This
is within reach of modern hardware but represents a board of
$1.4 \times 10^6 \times 1.4 \times 10^6 \approx 2 \times 10^{12}$
cells --- far larger than the standard $256 \times 256$ board.

\subsection{Comparison with biological estimates}

For context, estimates of the probability of spontaneous emergence of a
self-replicating RNA molecule range from $10^{-40}$ to $10^{-100}$ per
random nucleotide sequence of appropriate length. Our estimate of
$\sim 2^{-49} \approx 10^{-14.7}$ for the BRK variant reflects the
relative simplicity of the 6502 instruction set compared to
biochemistry: a 6502 replicator requires only $\sim$49 bits of specific
information, whereas even the simplest ribozyme replicators require
hundreds of specific nucleotides.

%----------------------------------------------------------------------
\section{Sensitivity Analysis}
\label{sec:sensitivity}

\subsection{Undocumented NOP counts}

The counts $n_1 = 7$, $n_2 = 13$, $n_3 = 7$ depend on the specific
6502 variant emulated. The NMOS 6502 has more undocumented NOPs than
the CMOS 65C02. The Sfotty emulator used in 6502life may implement
a different subset. The NOP correction factor is relatively insensitive
to these counts:

\begin{center}
\begin{tabular}{ccc}
\toprule
$(n_1, n_2, n_3)$ & $Q(1)$ & $F_{\mathrm{NOP}} = Q(1)^4$ \\
\midrule
$(1, 0, 0)$ & $256/255 = 1.0039$ & $1.0157$ \\
$(7, 0, 0)$ & $256/249 = 1.0281$ & $1.1174$ \\
$(7, 13, 7)$ & $256/229 = 1.1179$ & $1.5608$ \\
$(7, 20, 7)$ & $256/222 = 1.1532$ & $1.7695$ \\
\bottomrule
\end{tabular}
\end{center}

The NOP correction ranges from negligible ($+0.02$ bits) to modest
($+0.82$ bits), confirming that it is a minor effect.

\subsection{Functional but imperfect replicators}

Our analysis counts only ``perfect'' byte-level replicators that
produce an exact copy of page~0. A weaker definition might allow
replicators that copy enough code for the child to also replicate,
even if some non-code bytes are incorrect. This could significantly
increase the probability if the replicator were robust to background
byte variations, but the minimal replicators we study copy the entire
page and are therefore sensitive to the background (the child's page~0
is an exact copy of the parent's page~0, including any random garbage
in non-code bytes).

\subsection{Replicator as a substring}

Our analysis assumes the replicator starts at byte~0 (the entry point).
An alternative scenario: the replicator code exists somewhere in the
1024-byte cell (not necessarily at byte~0), and some sequence of
instructions at byte~0 eventually transfers control to it. This is
much harder to quantify, as it depends on the probability that random
code produces a jump to the correct offset. We conjecture this
increases $\Prep$ by at most a factor of $\sim 1024$ (one extra trial
per byte position), reducing $\Beff$ by at most $\sim 10$ bits. But
the probability that random prefix code successfully reaches the
replicator without crashing or looping is very low, likely offsetting
this gain.

%----------------------------------------------------------------------
\section{Dominance Time and Lottery Economics}
\label{sec:dominance}

\subsection{Replicator dominance time}

A replicator seeded at one cell spreads via the checkerboard pairing
scheme. Each half-pass, the replicator may copy to its paired neighbor
if it is the active cell in the pair (probability $\frac{1}{2}$) and
the quantum budget suffices for the copy loop (which it does for any
budget $\gtrsim 100$ cycles, essentially always).

The spreading wavefront advances at most 1 cell per half-pass in one
cardinal direction. The two half-passes in each round use orthogonal
tilings (horizontal and vertical), so the wavefront can advance 1 cell
horizontally and 1 cell vertically per round. Including the random
offset bits (which shift the tiling grid by 0 or 1), every cell
boundary is crossed by some pairing within $O(1)$ rounds.

For an $L \times L$ board, the replicator must traverse at most $L$
cells in each dimension. With a 50\% activation probability per pair,
the expected number of rounds to advance one cell is 2 (geometric
with $p = \frac{1}{2}$). The expected dominance time is therefore:
\[
D(L) \approx 2L \text{ rounds} = 4L \text{ half-passes (runPass calls).}
\]
Each half-pass gives each cell one quantum of $\mu = 4096$ mean cycles.
In cycles per cell:
\[
D_{\text{cyc}}(L) = D(L) \times \mu = 4L \times 4096 \approx 16{,}384 \, L.
\]
For $L = 32$: $D_{\text{cyc}} \approx 524{,}288$ cycles/cell $\approx
0.5 \times 10^6$. At $C = 200{,}000$ cyc/cell/s, dominance takes
$\approx 2.6$ seconds.

Note this is a lower bound on the mean; the actual distribution has a
tail from the random walk of the wavefront boundary. A more precise
analysis using first-passage times of a biased 2D random walk gives
$D(L) = \Theta(L)$ with moderate variance ($\sigma \sim \sqrt{L}$).

\subsection{Throughput: boards tested per unit time}

A player initializes a board with SHA-256 pseudorandom data and runs
the simulation to detect whether a viable replicator emerges. Detection
requires running long enough for a replicator (if present) to achieve
observable dominance---say, colonizing $> 50\%$ of cells.

At rate $C$ cycles/cell/s on an $L \times L$ board:
\begin{itemize}
\item Rounds per second: $C / (2\mu) \approx C / 8192$.
\item Time to dominance: $D(L) / (C / 2\mu) = 4L \cdot 2\mu / C
  = 8\mu L / C$ seconds.
\item Boards per second: $C / (8\mu L) \approx C / (32{,}768 \, L)$.
\end{itemize}

For $L = 32$, $C = 200{,}000$:
\begin{align*}
\text{Time per board} &= 8 \times 4096 \times 32 / 200{,}000
  \approx 5.2 \text{ seconds}, \\
\text{Boards per hour} &\approx 690.
\end{align*}

For $L = 16$, $C = 800{,}000$ (smaller board, faster per-cell rate
since $C \propto L^{-2}$ on fixed hardware):
\begin{align*}
\text{Time per board} &\approx 0.65 \text{ seconds}, \\
\text{Boards per hour} &\approx 5{,}500.
\end{align*}

\subsection{Lottery ticket probability}

Each cell in a randomly initialized board is an independent trial for
containing a replicator at position 0. With $L^2$ cells, the expected
number of replicators per board is:
\[
\mathbb{E}[\text{replicators per board}] = L^2 \cdot 2^{-\Beff}.
\]

For a player seeking at least one replicator, the expected number of
boards until success is:
\[
N_{\text{boards}} = \frac{1}{L^2 \cdot 2^{-\Beff}} = \frac{2^{\Beff}}{L^2}.
\]

\paragraph{Random initialization ($\Beff \approx 55$ bits).}
For $L = 32$: $N_{\text{boards}} = 2^{55} / 1024 = 2^{45} \approx
3.5 \times 10^{13}$. At 690 boards/hour, this requires $\sim 5 \times
10^{10}$ hours $\approx 5.8$ million years. \textbf{Infeasible.}

\paragraph{Zero initialization ($\Beff \approx 41$ bits).}
If RAM is zero-initialized and only the first few bytes are randomized
(not all 1024), the BRK variant has $\Beff \approx 41$:
$N_{\text{boards}} = 2^{41} / 1024 = 2^{31} \approx 2.1 \times 10^9$.
At 690 boards/hour: $\sim 3 \times 10^6$ hours $\approx 350$ years.
\textbf{Still infeasible for a single player, but conceivable for a
network.}

\paragraph{BRK-copy primitive ($\Beff \approx 12.4$ bits).}
If the VM's BRK-copy operation is enabled (operands 49--96 trigger
a built-in noisy cell copy), a ``replicator'' is simply \opcode{BRK}
followed by one of 48 valid operands. This is 2 constrained bytes
with 48/256 probability for the second: $P = 48/65536 \approx
2^{-10.4}$ per cell. For a $32 \times 32$ board:
$\mathbb{E} = 1024 \times 2^{-10.4} \approx 0.75$ replicators per
board. With $\approx 1.3$ boards, a player expects success.
\textbf{Trivially achievable}---but BRK-copy replicators are noisy
and may not sustain themselves against mutation.

\paragraph{Summary of lottery economics.}

\begin{center}
\begin{tabular}{lcccc}
\toprule
Scenario & $\Beff$ & $N_{\text{boards}}$ ($L\!=\!32$) &
  Time (690 brd/hr) & Feasible? \\
\midrule
JMP, random      & 65 & $3 \times 10^{17}$ & $5 \times 10^{10}$ yr & No \\
Branch, random   & 55 & $3 \times 10^{13}$ & $5 \times 10^{6}$ yr & No \\
BRK, zero-init   & 41 & $2 \times 10^{9}$  & 350 yr & Network \\
BRK-copy primitive & 10 & 1.3 & 7 sec & Yes \\
\bottomrule
\end{tabular}
\end{center}

The stark conclusion: programmatic self-replicators (byte-copy loops)
are astronomically unlikely to arise by chance. The VM's built-in
BRK-copy mechanism, if enabled, is the only plausible path to
spontaneous replication---but it produces noisy, fragile replicators
that must evolve error tolerance to persist.

\section{Conclusion}
\label{sec:conclusion}

\begin{center}
\begin{tabular}{lccc}
\toprule
& Strict $\Pbase$ & Relaxed $\Prep$ & $\Beff$ \\
\midrule
JMP (9 bytes) & $2^{-72}$ & $2^{-64.8}$ & 64.8 \\
Branch (8 bytes) & $2^{-64}$ & $2^{-55.2}$ & 55.2 \\
BRK, random init (8 bytes)\footnotemark & $2^{-64}$ & $2^{-56.8}$ & 56.8 \\
BRK, zero init (6 bytes) & $2^{-48}$ & $2^{-40.8}$ & 40.8 \\
\bottomrule
\end{tabular}
\end{center}

\footnotetext{BRK operand byte (byte 7) must be \$00 under random
initialization; under zero-init it is free.}

The minimal self-replicator in the 6502life VM has an effective
information content of approximately 41--57 bits, depending on the
assumed initial conditions. Spontaneous emergence on a standard
$256 \times 256$ board is negligible ($\sim 10^{-8}$ to $10^{-15}$
expected replicators). A board of side $\sim 10^6$ (requiring
$\sim$1.8 TiB) would be needed for spontaneous replication to be
expected under the most favorable assumptions.

The dominant relaxations are:
\begin{enumerate}
\item \textbf{Target cell flexibility} (byte~4): 48 valid neighbors
  instead of 1, saving $\approx 5.6$ bits.
\item \textbf{Bidirectional stepping} (byte~5): \opcode{INX} or
  \opcode{DEX}, saving 1 bit.
\item \textbf{NOP insertion tolerance}: saving $\approx 0.6$ bits
  (minor effect due to the low probability of NOP opcodes).
\end{enumerate}

The remaining bytes (load opcode, base addresses, store opcode, loop
mechanism) are tightly constrained with no or very few alternatives.

\begin{thebibliography}{9}
\bibitem{6502life}
6502life.com: A virtual 256$\times$256 grid of interconnected 6502 CPUs
simulating cellular automata.
\url{https://6502life.com}, 2024--2026.
\end{thebibliography}

\end{document}
